\documentclass[12pt,a4paper]{article}
\usepackage[utf8]{inputenc}
\usepackage[english,greek]{babel} % for greek
%\usepackage[T1]{fontenc}
\usepackage{amsmath}
\usepackage[svgnames]{xcolor}
\usepackage{amsfonts}
\usepackage{amssymb}
\usepackage{graphicx}
\usepackage{hyperref}
\usepackage{placeins}
\usepackage[left=1.50cm, right=1.50cm, top=3.00cm, bottom=3.00cm]{geometry}
%\author{Stratos Pateloudis}
\title{Καιρός και ελιές}
\begin{document}
\maketitle
Σκοπός αυτής της ανάλυσης είναι η εξερεύνηση πιθανών συσχετίσεων μεταξύ των καιρικών συθηκών και ποσοτικών χαρακτηριστικών των ελιών. Για το λόγο αυτό συλλέχτηκαν δεδομένα απο το 2009 που αφορούν τη θερμοκρασία, ύψος βροχοπτώσεων και χιονοπτώσεων, πιέσεις, ανέμους, ηλιοφάνεια, συννεφοκάλυψη, κλπ. Αξίζει να σημειωθεί ότι τα στατιστικά στοιχεία δεν επαρκούν για την εξαγωγή ασφαλών συμπερασμάτων αλλά αποτελούν ένα δείγμα και ίσως δίνουν μια γενικη τάση.	
\section{Γενική ανάλυση καιρικών συνθηκών}
Το μικροκλίμα της Νέας Ζίχνης έχει τα εξής ετήσια χαρακτηριστικά: η μέση θερμοκρασία κυμαίνεται στους $15.87$ βαθμούς Κελσίου, η αθροιστική μέση βροχόπτωση βρίσκεται στα $728.35$ χιλιοστά, η χιονόπτωση στα $0.6$ εκατοστά, η ηλιοφάνεια στις $244$ ώρες, η συνεφοκάλυψη στο $30\%$, και οι άνεμοι στα $6.9$ χλμ/ώρα. Παραδοσιακά η περίοδος των βροχοπτώσεων είναι στο τέταρτο τρίμηνο του έτους με σημαντική εξαίρεση την χρονιά του 2023.% όπως φαίνεται στο σχήμα \ref{fig:quarterly_seasonal}.

H μηνιαία βροχόπτωση ανά έτος φαίνεται στο σχήμα \ref{fig:rain_seasonal}. Παρατηρείται ότι τα μεγαλύτερα ύψη ξεκινούν το φθινόπωρο όπως προείπαμε (απο Σεπτέμβριο) μέχρι και τον χειμώνα, με εξαίρεση τον Μάιο του 2023.

\section{Συχέτιση παραγωγής και καιρού}
Η παραγωγή φαίνεται πως συσχετίζεται με ένα γραμμικό μοντέλο που εξαρτάται απο τις ακόλουθες καιρικές μεταβλητές: Μέγιστοι άνεμοι τον Μάρτιο, Μέσες θερμοκρασίες Απρίλιο και Σεπέμβρη, Μέγιστη θερμοκρασία Ιούνιο, Συννεφοκάλυψη τον Αύγουστο. Απο το μοντέλο φαίνεται ότι θετικά συνεισφέρουν οι μεταβλητές: Μέση θερμοκρασία τον Απρίλιο και συννεφοκάλυψη τον Αύγουστο ενώ αρνητικά οι υπόλοιπες. 

Γραφικές αναπαραστάσεις αυτών των καιρικών μεταβλητών καθώς και μέσοι όροι δίνονται στα σχήματα παρακάτω. 
%\begin{equation*}
%	yield = -251.66*Maxwind_3 + 742.04*Avgtemp_4 - 265.35 *Maxtemp_6 + 202.93*Clouds_8 - 207.52*Avgtemp_9  
%\end{equation*}

\begin{figure}[ht!]
	\centering 
	\includegraphics[scale=0.67]{Rain_seasonal}
	\caption{Ύψος βροχοπτώσεων ανα μήνα και έτος.}
	\label{fig:rain_seasonal}
\end{figure}

\begin{figure}[ht!]
	\centering 
	\includegraphics[scale=0.6]{Avg_temp}
	\caption{Μέσες θερμοκρασίες ανά έτος. Η κουκίδα αντιπροσωπεύει τον μέσο όρο τον συγκεκριμένο μήνα απο όλες τις χρονιές.}
	\label{fig:mean_temp_history}
\end{figure}
\begin{figure}[ht!]
	\centering 
	\includegraphics[scale=0.6]{max_temp}
	\caption{Μέγιστες Θερμοκρασίες. Η κουκίδα αντιπροσωπεύει τον μέσο όρο τον συγκεκριμένο μήνα απο όλες τις χρονιές.}
	\label{fig:Sum_Rain_history}
\end{figure}

\begin{figure}[ht!]
	\centering 
	\includegraphics[scale=0.6]{clouds}
	\caption{Συννεφοκάλυψη. Η κουκίδα αντιπροσωπεύει τον μέσο όρο τον συγκεκριμένο μήνα απο όλες τις χρονιές}
	\label{fig:Sum_Rain_history}
\end{figure}


\section{Πρώτα συμπεράσματα}
Το μοντέλο εξηγεί την παραγωγή κάθε χρονιάς με ακρίβεια 200 κιλών, λαμβάνοντας υπόψιν τις παραπάνω καιρικές μεταβλητές.  
\begin{figure}[ht!]
	\centering 
	\includegraphics[scale=0.6]{predictions}
	\caption{Προβλέψεις (μπλέ) και πραγματικές παραγωγές (πορτοκαλί) με βάση μόνο τις καιρικές συνθήκες.}
	\label{fig:Sum_Rain_history}
\end{figure}
 
	
\end{document}